% SYNTHETIC EVALUATION DRAFT — written for the retrieval eval in evals/.
% Not a real paper; contains no experimental results.
\documentclass{article}
\title{Mitigating Over-Smoothing and Over-Squashing in Deep Message-Passing Graph Neural Networks}
\begin{document}
\maketitle

\begin{abstract}
Message-passing graph neural networks (GNNs) learn node representations by
iteratively aggregating features from neighbouring nodes. Stacking many layers
should enlarge the receptive field, but in practice deep GNNs suffer from
over-smoothing, where node embeddings become indistinguishable, and
over-squashing, where information from distant nodes is compressed through
structural bottlenecks. This draft proposes a curvature-guided graph rewiring
step combined with an initial-residual connection, and outlines an evaluation
on node classification and graph-level property prediction.
\end{abstract}

\section{Introduction}
Graphs model relational data such as citation networks, social networks,
molecules and knowledge graphs. Graph neural networks generalize convolution to
irregular graph structure: in each layer, every node aggregates messages from
its neighbours and updates its hidden state. Graph convolutional networks
(GCN), GraphSAGE, graph attention networks (GAT) and graph isomorphism
networks (GIN) are widely used instances of this message-passing paradigm.

Two failure modes limit depth. Over-smoothing arises because repeated
neighbourhood averaging behaves like a low-pass filter on the graph Laplacian,
driving node features toward a common subspace. Over-squashing arises because
the number of nodes in a $k$-hop neighbourhood grows exponentially, while the
messages passed across a bottleneck edge have fixed dimension. Long-range
dependencies, which are common in molecular graphs, are therefore poorly
captured.

Our contributions are: (1) a rewiring procedure that adds edges across
negatively curved bottleneck edges, measured by a discrete Ricci curvature,
and removes edges in highly positively curved regions; (2) an initial-residual
message-passing layer that mixes each layer's output with the input features
to slow over-smoothing; and (3) a study of the trade-off between the two
phenomena as depth increases.

\section{Related Work}
Over-smoothing has been analysed through the spectrum of the normalized
adjacency matrix, and mitigated with residual connections, jumping knowledge
networks, PairNorm, DropEdge and APPNP-style personalized PageRank
propagation. Over-squashing was linked to the Jacobian sensitivity between
distant nodes and to negative Ricci curvature, motivating graph rewiring
methods such as SDRF and spectral approaches based on the Cheeger constant.
Graph transformers sidestep both issues with global attention but at quadratic
cost in the number of nodes. Expressive power of message passing is bounded by
the Weisfeiler--Lehman graph isomorphism test.

\section{Method}
\subsection{Curvature-Guided Rewiring}
We compute the balanced Forman curvature of every edge once before training.
For the most negatively curved edges, we add a supporting edge between
neighbours of the two endpoints that maximally increases curvature. To keep
the edge count roughly constant and avoid creating new over-smoothing, we
remove edges whose curvature exceeds an upper threshold.

\subsection{Initial-Residual Message Passing}
Each layer computes $H^{(l+1)} = \sigma\big((1-\beta)\hat{A}H^{(l)} +
\beta H^{(0)}\big)W^{(l)}$, where $\hat{A}$ is the normalized adjacency of the
rewired graph and $\beta$ controls how much of the input signal is preserved.

\section{Experimental Plan}
We plan to evaluate on homophilous citation graphs, heterophilous web graphs
and molecular property prediction datasets that require long-range
interactions. Baselines include GCN, GAT, GIN, GCNII, DropEdge, SDRF rewiring
and a graph transformer. We will vary depth from 2 to 64 layers and measure
accuracy together with the Dirichlet energy of node embeddings as an
over-smoothing diagnostic.

\section{Limitations}
Curvature computation scales with the number of edges times node degree,
which is costly on very large graphs. Rewiring changes the input topology,
which may be undesirable when edges carry semantic meaning.

\end{document}
